%%%%%%%%%%%%%%%%%%%%%%%%%%%%%%%%%%%%%%%%%%%%%%%%%%%%%%%%%%%%%%%%%%%%%%%%%%%%%%%
%%  Projeto detalhado para o Comitê de Ética (gera o etica.pdf).
%%  Siga as Resoluções CNS 466/2012 e CNS 510/2016. Substitua os lembretes
%%  \todo[inline] pelo seu texto e remova-os no final.
%%%%%%%%%%%%%%%%%%%%%%%%%%%%%%%%%%%%%%%%%%%%%%%%%%%%%%%%%%%%%%%%%%%%%%%%%%%%%%%
\section{Resumo}
\todo[inline]{Resuma em um parágrafo o problema, o objetivo, quem participará e o que será feito com os participantes.}

\section{Introdução e justificativa}
\todo[inline]{Contextualize o tema, apresente a lacuna que a pesquisa preenche e justifique por que é necessário envolver seres humanos.}

\section{Hipótese}
\todo[inline]{Descreva a hipótese da pesquisa. Se o estudo for exploratório, justifique a ausência de hipótese.}

\section{Objetivos}
\subsection{Objetivo primário}
\todo[inline]{Um único objetivo, iniciado por verbo no infinitivo.}
\subsection{Objetivos secundários}
\todo[inline]{Objetivos específicos que contribuem para o objetivo primário.}

\section{Metodologia}
\subsection{Desenho do estudo}
\todo[inline]{Tipo de estudo (por exemplo, experimental, observacional, qualitativo, survey).}
\subsection{Participantes}
\todo[inline]{Como os participantes serão recrutados e quantos serão (justifique o número). Local: \eticaLocalColeta. Público: \eticaPublicoAlvo. Total previsto: \eticaNumeroParticipantes.}
\subsubsection*{Critérios de inclusão}
\todo[inline]{Quem pode participar.}
\subsubsection*{Critérios de exclusão}
\todo[inline]{Quem não pode participar e por quê.}
\subsection{Procedimentos}
\todo[inline]{Descreva, em ordem, o que cada participante fará, a duração de cada etapa e os instrumentos usados (questionários, tarefas, equipamentos). Anexe os instrumentos ao processo na Plataforma Brasil.}
\subsection{Análise dos dados}
\todo[inline]{Como os dados serão analisados.}

\section{Riscos e benefícios}
\subsection{Riscos}
\todo[inline]{Todo estudo com seres humanos tem algum risco (cansaço, constrangimento, quebra de sigilo, desconforto). Liste cada um com a probabilidade e a gravidade.}
\subsection{Medidas para minimizar os riscos}
\todo[inline]{Para cada risco, diga como será reduzido e o que será feito se ele ocorrer (por exemplo, interromper a participação sem prejuízo).}
\subsection{Benefícios}
\todo[inline]{Benefícios diretos aos participantes (se houver) e benefícios para a sociedade e a ciência.}

\section{Desfechos}
\subsection{Desfecho primário}
\todo[inline]{O principal resultado que será medido.}
\subsection{Desfechos secundários}
\todo[inline]{Outros resultados medidos.}

\section{Confidencialidade e armazenamento dos dados}
\todo[inline]{Como a identidade dos participantes será protegida (códigos no lugar de nomes, acesso restrito), onde os dados ficarão guardados, por quanto tempo (mínimo de 5 anos) e quem terá acesso. Considere a LGPD (Lei 13.709/2018).}

\section{Ressarcimento e indenização}
\todo[inline]{Informe se haverá ressarcimento de despesas (transporte, alimentação) e a garantia de indenização em caso de dano decorrente da pesquisa.}

\section{Cronograma}
\todo[inline]{A coleta de dados só pode começar após a aprovação do CEP. Liste as atividades e o período de cada uma.}
\begin{longtable}{@{\extracolsep{\fill}}p{0.6\textwidth} l@{}}
  \toprule
  \textbf{Atividade} & \textbf{Período} \\
  \midrule
  \endhead
  Submissão à Plataforma Brasil e aprovação pelo CEP & 01/2027 a 02/2027 \\ % Alterar
  Recrutamento e coleta de dados & \eticaPeriodoColeta \\
  Análise dos dados & 06/2027 \\ % Alterar, adicionar ou remover
  Redação e entrega dos resultados & 07/2027 \\
  \bottomrule
\end{longtable}

\section{Orçamento}
\todo[inline]{O orçamento vem dos arquivos da pasta recursos/ (materiais, serviços e outros). Financiamento: \eticaFinanciamento.}
\tabelasRecursos{\subsection*}
\resumoRecursos{\subsection*}
